\documentclass[11pt]{amsart}
\usepackage[a4paper,margin=30mm]{geometry}
\usepackage{amsmath,amssymb,amsthm,mathtools}
\usepackage{microtype}
\usepackage{enumitem}
\usepackage{hyperref}

\newtheorem{theorem}{Theorem}[section]
\newtheorem{lemma}[theorem]{Lemma}
\newtheorem{corollary}[theorem]{Corollary}
\theoremstyle{remark}
\newtheorem{remark}[theorem]{Remark}

\newcommand{\cyc}{\sum_{\mathrm{cyc}}}

\title{A Proof of Liu's Conjecture on the Fundamental Triangle Inequality}
\author[Ts. Batbold]{Tserendorj Batbold}
\address{
Department of Mathematics\\
National University of Mongolia\\
Ulaanbaatar 14201\\
Mongolia}
\email{tsbatbold@hotmail.com}

\begin{document}


\keywords{triangle inequality; circumradius; inradius; Ravi substitution; Liu's conjecture}

\subjclass[2020]{51M16, 26D15.}
\begin{abstract}
Let $a,b,c$ be the side lengths of a triangle, and let $R$ and $r$
denote its circumradius and inradius, respectively. We prove a conjecture
of Liu stating that
\[
 \cyc \left(\frac{a(b+c-a)}{bc}\right)^k
 \ge 2+\left(\frac{2r}{R}\right)^k, \qquad k>1,
\]
with the reverse inequality for $0<k<1$. The proof reduces the problem
to three positive variables with fixed sum and product. We also determine
the equality cases.
\end{abstract}

\maketitle

\section{Introduction}
Throughout the paper, $a,b,c$ denote the side lengths of a nondegenerate triangle $ABC$, while $s,R,r$ denote its semiperimeter, circumradius, and inradius, respectively. We use the standard notation for triangle elements as in Mitrinovi\'c, Pe\v{c}ari\'c and Volenec \cite{Mitrinovic1989}.

In \cite{Liu2022}, Liu studied the fundamental triangle inequality and Gerretsen's double inequality and, among other results, proposed the following conjecture (Conjecture 5.1 in that paper): for every real $k>1$,
\begin{equation}\label{eq:liu}
 \cyc \frac{a^k(b+c-a)^k}{(bc)^k}
 \ge 2+\left(\frac{2r}{R}\right)^k,
\end{equation}
and the direction of \eqref{eq:liu} is reversed when $0<k<1$. A subsequent correction to the original paper concerns the numbering of its references \cite{LiuCorrection2022}. For $k=1$, Liu observed the identity
\begin{equation}\label{eq:k1}
 \cyc \frac{a(b+c-a)}{bc}=2+\frac{2r}{R}.
\end{equation}

Our proof is based on two elementary identities for the three terms in \eqref{eq:k1}. Their sum and product reduce the conjecture to an algebraic problem for three positive variables subject to fixed sum and product. We first establish the auxiliary result needed in the proof and then apply it to the triangle.

\section{Proof of the conjecture}
We begin with the following auxiliary result.

\begin{lemma}\label{lem:power}
Let $0<q\le 1$ and let $x,y,z>0$ satisfy
\begin{equation}\label{eq:constraints}
 x+y+z=2+q,\qquad xyz=q.
\end{equation}
Then
\begin{equation}\label{eq:power}
 x^k+y^k+z^k
 \begin{cases}
 \ge 2+q^k, & k>1,\\[2mm]
 \le 2+q^k, & 0<k<1.
 \end{cases}
\end{equation}
If $k\ne1$, equality holds if and only if $(x,y,z)$ is a permutation of $(q,1,1)$.
\end{lemma}

\begin{proof}
If $q=1$, then $x+y+z=3$ and $xyz=1$. Equality in AM--GM gives $x=y=z=1$, and the assertion follows. Hence assume $0<q<1$.

Consider
\[
 \mathcal S_q=\{(x,y,z)\in(0,\infty)^3:x+y+z=2+q,\ xyz=q\}.
\]
This set is nonempty, since $(q,1,1)\in\mathcal S_q$, and it is compact. Indeed, every coordinate is at most $2+q$. Moreover,
\[
 yz\le \left(\frac{y+z}{2}\right)^2\le\frac{(2+q)^2}{4},
\]
so
\[
 x=\frac{q}{yz}\ge\frac{4q}{(2+q)^2},
\]
and cyclically the same positive lower bound holds for $y$ and $z$. Hence $\mathcal S_q$ is a closed subset of a compact box. Thus the continuous function
\[
 F(x,y,z)=x^k+y^k+z^k
\]
attains its minimum and maximum on $\mathcal S_q$.

We next verify that the constraint gradients are independent on $\mathcal S_q$. If
\[
 (1,1,1)\quad\text{and}\quad(yz,zx,xy)
\]
were linearly dependent, positivity would imply $x=y=z$. Then \eqref{eq:constraints} would yield
\[
 \left(\frac{2+q}{3}\right)^3=q,
\]
or
\[
 (2+q)^3-27q=(q-1)^2(q+8)=0,
\]
contrary to $0<q<1$. Hence the Lagrange multiplier conditions apply at every extremum of $F$ on $\mathcal S_q$.

At an extremum there exist real $\lambda,\mu$ such that
\[
 kx^{k-1}=\lambda+\mu yz,
\]
and cyclically. Since $xyz=q$, multiplication by the corresponding variable gives
\[
 kx^k-\lambda x-\mu q=0,
\]
with the corresponding equations for $y$ and $z$. Hence $x,y,z$ are positive zeros of
\[
 \phi(t)=kt^k-\lambda t-\mu q.
\]
For $t>0$,
\[
 \phi''(t)=k^2(k-1)t^{k-2},
\]
which has a fixed nonzero sign when $k\ne1$. Thus $\phi$ is strictly convex for $k>1$ and strictly concave for $0<k<1$. A strictly convex function, as well as a strictly concave function, has at most two distinct zeros on an interval: if three zeros existed, strict convexity or strict concavity applied to the middle zero and the two outer zeros would give a strict inequality against $0$. Therefore at least two of $x,y,z$ are equal at every extremum.

By symmetry, write $y=z=t$. From \eqref{eq:constraints},
\[
 x=\frac{q}{t^2},\qquad \frac{q}{t^2}+2t=2+q.
\]
Hence
\[
 2t^3-(2+q)t^2+q=0,
\]
which factors as
\begin{equation}\label{eq:factor}
 (t-1)(2t^2-qt-q)=0.
\end{equation}

If $t=1$, then $(x,y,z)=(q,1,1)$ and
\begin{equation}\label{eq:F1}
 F_1=2+q^k.
\end{equation}

For the other positive solution of \eqref{eq:factor},
\[
 q=\frac{2t^2}{1+t},\qquad x=\frac{2}{1+t}.
\]
Since $0<q<1$, this positive solution satisfies $0<t<1$. The corresponding value of $F$ is
\[
 F_2=\left(\frac{2}{1+t}\right)^k+2t^k.
\]
Using $q=2t^2/(1+t)$, we obtain
\begin{align*}
 F_2-F_1
 &=\left(\frac{2}{1+t}\right)^k+2t^k-2
   -\left(\frac{2t^2}{1+t}\right)^k\\
 &=(1-t^k)\left[\left(\frac{2}{1+t}\right)^k(1+t^k)-2\right].
\end{align*}
For $k>1$, strict convexity of $u\mapsto u^k$ gives
\[
 \frac{1+t^k}{2}>\left(\frac{1+t}{2}\right)^k,
\]
so $F_2>F_1$. Thus the global minimum is $F_1$, proving the first inequality in \eqref{eq:power}. For $0<k<1$, strict concavity reverses the last inequality and yields $F_2<F_1$; hence the global maximum is $F_1$, proving the second inequality.

The strictness of the comparison for $0<t<1$ shows that equality can occur only at a permutation of $(q,1,1)$. Conversely, every permutation of \((q,1,1)\) satisfies the constraints and gives equality.
\end{proof}

\begin{theorem}\label{thm:main}
Let $ABC$ be a nondegenerate triangle with side lengths $a,b,c$, circumradius $R$, and inradius $r$. If $k>1$, then
\begin{equation}\label{eq:mainplus}
 \cyc \left(\frac{a(b+c-a)}{bc}\right)^k
 \ge 2+\left(\frac{2r}{R}\right)^k.
\end{equation}
If $0<k<1$, then
\begin{equation}\label{eq:mainminus}
 \cyc \left(\frac{a(b+c-a)}{bc}\right)^k
 \le 2+\left(\frac{2r}{R}\right)^k.
\end{equation}
For $k\ne1$, equality holds if and only if $ABC$ is isosceles.
\end{theorem}

\begin{proof}
Use the Ravi substitution
\[
 a=v+w,\qquad b=w+u,\qquad c=u+v,
\]
where $u=s-a$, $v=s-b$, $w=s-c$ are positive. Set
\[
 X=\frac{a(b+c-a)}{bc},\qquad
 Y=\frac{b(c+a-b)}{ca},\qquad
 Z=\frac{c(a+b-c)}{ab}.
\]
Then
\[
 X=\frac{2u(v+w)}{(u+v)(u+w)},\quad
 Y=\frac{2v(w+u)}{(v+w)(u+v)},\quad
 Z=\frac{2w(u+v)}{(u+w)(v+w)}.
\]
A direct multiplication gives
\begin{equation}\label{eq:prodXYZ}
 XYZ=\frac{8uvw}{(u+v)(v+w)(w+u)}.
\end{equation}
Also, direct simplification gives
\begin{equation}\label{eq:sumXYZ}
 X+Y+Z=2+\frac{8uvw}{(u+v)(v+w)(w+u)}.
\end{equation}
The expression on the right-hand side of \eqref{eq:prodXYZ} has a standard geometric form. Heron's formula gives
\[
 \Delta^2=(u+v+w)uvw,
\]
while
\[
 r=\frac{\Delta}{u+v+w},\qquad
 R=\frac{(u+v)(v+w)(w+u)}{4\Delta}.
\]
Therefore,
\begin{equation}\label{eq:q}
 \frac{2r}{R}=\frac{8uvw}{(u+v)(v+w)(w+u)}.
\end{equation}
Combining \eqref{eq:prodXYZ}--\eqref{eq:q},
\[
 X+Y+Z=2+q,\qquad XYZ=q,
 \qquad q:=\frac{2r}{R}.
\]
By Euler's inequality $R\ge2r$, we have $0<q\le1$. Applying Lemma~\ref{lem:power} to $X,Y,Z$ gives \eqref{eq:mainplus} and \eqref{eq:mainminus}.

We now determine the equality case. By Lemma~\ref{lem:power}, for $k\ne1$ equality is equivalent to $(X,Y,Z)$ being a permutation of $(q,1,1)$. Suppose $Y=Z=1$. From the formulas above,
\[
 Y=1\iff (v-w)(u-v)=0,
\]
and
\[
 Z=1\iff (w-v)(u-w)=0.
\]
If $v\ne w$, these two relations would force $u=v=w$, a contradiction. Hence $v=w$, equivalently $b=c$. The other permutations give $c=a$ or $a=b$. Conversely, if two sides are equal, two of $X,Y,Z$ equal $1$ and the third equals $XYZ=q$, so equality follows. Thus equality holds exactly for isosceles triangles.
\end{proof}

\begin{remark}
At $k=1$, the assertion becomes the identity
\[
 \cyc \frac{a(b+c-a)}{bc}=2+\frac{2r}{R},
\]
which is exactly \eqref{eq:sumXYZ} together with \eqref{eq:q}.
\end{remark}

\section{Concluding remarks}
The identities obtained from Ravi's substitution show that the triangle problem is governed by the two relations
\[
 X+Y+Z=2+q,\qquad XYZ=q,\qquad q=\frac{2r}{R}\in(0,1].
\]
Lemma~\ref{lem:power} then gives the required inequalities for both ranges of $k$. Its strict form also yields the equality characterization: for $k\ne1$, equality occurs exactly for isosceles triangles. At $k=1$ the inequality reduces to the identity \eqref{eq:k1}. The same sum-product reduction may be useful for related power inequalities in triangle geometry.

\begin{thebibliography}{9}

\bibitem{Liu2022}
J. Liu,
\emph{On the fundamental triangle inequality and Gerretsen's double inequality},
J. Geom. \textbf{113} (2022), Article 19, 22 pp.
\href{https://doi.org/10.1007/s00022-022-00630-w}{doi:10.1007/s00022-022-00630-w}.

\bibitem{LiuCorrection2022}
J. Liu,
\emph{Correction to: On the fundamental triangle inequality and Gerretsen's double inequality},
J. Geom. \textbf{113} (2022), Article 27.
\href{https://doi.org/10.1007/s00022-022-00639-1}{doi:10.1007/s00022-022-00639-1}.

\bibitem{Mitrinovic1989}
D. S. Mitrinovi\'c, J. E. Pe\v{c}ari\'c and V. Volenec,
\emph{Recent Advances in Geometric Inequalities},
Kluwer Academic Publishers, Dordrecht, 1989.

\end{thebibliography}

\end{document}
